\subsection{Sprint 0}

O projeto tecnicamente teria iniciado após a primeira reunião com as
stakeholders do MP, que apresentaram sua proposta sobre um aplicativo de ajuda a
vítimas de crime. Porém, devido a possível confusão pelas duas partes (nossa e do
cliente), não foi possível determinar os requisitos e nem o escopo do projeto.

Tendo em vista este pequeno atraso, os AGES IV do projeto marcaram outra
reunião na mesma semana para definir de fato o que nos faltava para iniciar os
trabalhos. Após esta segunda reunião as primeiras tarefas já foram divididas, entre
escolha de tecnologias e a arquitetura para os AGES III, um quadro \textit{kanban} para
organizar as tasks e o design de mockups no figma das telas do app para os AGES I
e II.

Como AGES I e sem conhecimento de nenhuma tecnologia usada no projeto,
nesta \textit{sprint} acabei apenas observando e estudando o que teríamos que utilizar para
realização do produto. Em conjunto com outros AGES I desenvolvi algumas telas no
figma, porém meu foco foi maior em estudar o framework Spring Boot para o \textit{backend}
e React Native para o \textit{frontend}
\\
\\
\fbox{%
    \parbox{1\linewidth}{%
        \setlength{\parindent}{15pt}
            \itshape Minha experiência inicial de AGES I deixou muitíssimo a desejar. A não ser que o aluno em sua
            primeira AGES já  tenha trabalhado na área, é de extrema dificuldade que ele
            consiga exercer alguma função dentro do projeto.
    
            Há várias tecnologias que não são abordadas em nenhuma disciplina
            específica, e, em grande parte, espera-se que o aluno já tenha conhecimento e/ou
            experiência em seu uso. A meu ver isso tudo impacta para uma vivência que não
            agrega muito no currículo dos alunos.
    
            Não foi de ajuda também, o fato que o início do projeto atrasou pela confusão
            das clientes, além da cultura forte de herói no time.
    }%
}
